\chapter{Meta y Objetivos}

\section{Hipótesis}

El diseño y la integración de una estación inteligente de control ESD de bajo
costo permitirán transformar la verificación manual de pulseras antiestáticas en
un evento trazable y auditable. Para ello, el sistema integrará un probador
comercial HZR-171, mecanismos de identificación de usuarios, una interfaz
gráfica táctil y almacenamiento local de los registros generados durante cada
prueba.

El prototipo generará registros completos en al menos el 95\% de los ciclos de
prueba ejecutados. Cada registro deberá incluir el nombre y el identificador del
usuario, la fecha, la hora, el resultado ESD, la temperatura y la humedad relativa.
Además, el costo estimado de materiales se mantendrá por debajo de USD~350.

\section{Objetivos}

\subsection{Objetivo general}

Diseñar, implementar y validar un prototipo de estación inteligente de bajo
costo para el control y la trazabilidad de verificaciones ESD en laboratorios de
electrónica, mediante la integración de un probador de pulseras antiestáticas,
mecanismos de identificación de usuarios, una interfaz gráfica y almacenamiento
local de registros.

\subsection{Objetivos específicos}

\begin{itemize}

\item Diseñar la arquitectura electrónica, de comunicación, alimentación e
integración física de la estación, con base en los requerimientos
funcionales y no funcionales del sistema.

\item Diseñar e implementar los subsistemas de identificación de usuarios y
adquisición del resultado de la prueba ESD, mediante la integración del lector
de código de barras, el reconocimiento facial complementario y las señales del
probador HZR-171.

\item Diseñar e implementar la interfaz gráfica que guíe al usuario durante el
ciclo de prueba y el sistema de almacenamiento de los resultados y generación
de reportes históricos trazables.

\item Validar el desempeño funcional del prototipo mediante pruebas individuales
de los subsistemas y pruebas de integración del sistema completo, estas últimas
realizadas en un laboratorio de electrónica, utilizando los indicadores de
trazabilidad e integración funcional definidos para el proyecto.

\end{itemize}
